\clearpage
\section{A continuous Hecke limit and its complete category}
\label{sec:f1-limit-categories}

The positive Hecke parameter gives a genuine continuous path from arithmetic
flag algebras to the symmetric-group fibre at $q=1$.  The construction below
keeps one parameter throughout: tensor products of section algebras are
balanced over $C(I)$, whereas tensor products inside one fibre are ordinary
spatial tensor products.  Completing the parabolic objects then produces all
finite self-adjoint matrix corners, including off-diagonal maps between
summands that happen to have the same total degree.

\subsection{Continuous Hecke sections and parabolic corners}

\begin{definition}[Positive Hecke algebra and normalized regular frame]
Fix $q>0$.  Let $H_0(q)=H_1(q)=\mathbb C$, and for $n\geq2$ let $H_n(q)$
be generated by self-adjoint $T_1,\ldots,T_{n-1}$ subject to
\[
 (T_i-q)(T_i+1)=0,\qquad
 T_iT_{i+1}T_i=T_{i+1}T_iT_{i+1},\qquad
 T_iT_j=T_jT_i\quad(|i-j|>1).
\]
For $w\in S_n$, write $T_w$ for a reduced-word product and
$\ell(w)$ for its Coxeter length.  The coefficient trace is
$\tau_{n,q}(T_w)=\delta_{w,1}$.  Put
\[
 b_w(q)=q^{-\ell(w)/2}T_w(q).
\]
The vectors $b_w(q)$ are identified with the standard basis $\delta_w$ of
$K_n=\ell^2(S_n)$, and $B_w(q)$ denotes left multiplication by $b_w(q)$
on this fixed Hilbert space.
\end{definition}
\begin{scope}
The embedding in the fixed regular Hilbert space is named data.  It does not
make the Hecke generators independent of $q$, and no positivity is asserted
for a non-real or non-positive parameter.
\end{scope}
\provenance{D1101, D1201}{theory/sidequests/f1-limit/continuous-corners.md}{theory/checks/f1\_limit\_check.py (L1); theory/checks/f1\_operational\_check.py (H1,H2)}{theory/verdicts/f1-limit-adjudication.md}

\begin{definition}[Continuous positive Hecke section algebra]
For a nonempty compact interval $I\subset(0,\infty)$, define
\[
 H_n^{\mathrm{cts}}(I)
 =\left\{q\longmapsto\sum_{w\in S_n}f_w(q)B_w(q):
                   f_w\in C(I)\right\}
 \subset C\bigl(I,\operatorname{End}(K_n)\bigr).
\]
Operations are pointwise and the norm is the supremum operator norm.  Its
$C(I)$-valued trace is
\[
 \tau_n^{\mathrm{cts}}(x)(q)=\tau_{n,q}(x(q)).
\]
\end{definition}
\begin{scope}
This is a continuous section algebra over a compact positive interval.  It
is not a generic algebra over a Laurent-polynomial coefficient ring.
\end{scope}
\provenance{D1202}{theory/sidequests/f1-limit/continuous-corners.md}{finite-fibre probes L1,H1,H2; closedness is proved analytically}{theory/verdicts/f1-limit-adjudication.md}

\begin{theorem}[Continuous regular Hecke algebras]
{\normalfont\statusProved}\quad
For every nonempty compact $I\subset(0,\infty)$ and every $n\geq0$,
$H_n^{\mathrm{cts}}(I)$ is a unital C*-subalgebra of
$C(I,\operatorname{End}(K_n))$.  Evaluation at $q_0\in I$ is a surjective
*-homomorphism onto the faithful regular copy of $H_n(q_0)$, and
$\tau_n^{\mathrm{cts}}$ is continuous, tracial, positive and faithful.
\end{theorem}
\begin{scope}
The assertion is uniform in the rank and on every compact positive interval.
The finite checks test fibres; they do not prove norm-closedness or
continuity.
\end{scope}
\provenance{F1-LIM-CTS}{theory/sidequests/f1-limit/continuous-corners.md}{theory/checks/f1\_limit\_check.py (L1); theory/checks/f1\_operational\_check.py (H1,H2)}{theory/verdicts/f1-limit-adjudication.md}
\begin{proof}
On a length pair $w,s_iw$ with $\ell(s_iw)=\ell(w)+1$, left multiplication
by $T_i$ has matrix
\[
 \begin{pmatrix}0&\sqrt q\\ \sqrt q&q-1\end{pmatrix}
\]
in the normalized basis.  Hence all $B_w(q)$ vary continuously.  The Hecke
relations show that their $C(I)$-span is a unital *-algebra.  If a sequence
of sections converges uniformly, applying it to $\delta_1$ recovers every
coefficient function uniformly; the limit therefore remains in the displayed
span.  This proves closedness.  Constant coefficient functions prove
surjectivity of every evaluation.  Finally,
\[
 \tau_{n,q}(x^*x)=\sum_w |c_w|^2q^{\ell(w)}
 \quad\text{for }x=\sum_wc_wT_w,
\]
so fibre traces, and consequently the section trace, are faithful and
positive.
\end{proof}

\begin{definition}[Continuous parabolic corner category]
For a composition $\alpha=(\alpha_1,\ldots,\alpha_r)$ of $n$, let
$W_\alpha\leq S_n$ preserve its consecutive blocks and set
\[
 P_\alpha(q)=\sum_{w\in W_\alpha}q^{\ell(w)},\qquad
 e_\alpha(q)=P_\alpha(q)^{-1}\sum_{w\in W_\alpha}T_w(q).
\]

The objects of $\Gamma_I^{\mathrm{cts}}$ are compositions, including the
empty composition of zero, and
\[
 \operatorname{Hom}(\alpha,\beta)=
 \begin{cases}
 e_\beta H_n^{\mathrm{cts}}(I)e_\alpha,&|\alpha|=|\beta|=n,\\
 0,&|\alpha|\neq|\beta|.
 \end{cases}
\]
Composition, identity and dagger are multiplication, $e_\alpha$ and matrix
adjoint.  Ordered tensor concatenates compositions.  On section algebras its
domain is the balanced product
\[
 H_m^{\mathrm{cts}}(I)\otimes_{C(I)}H_n^{\mathrm{cts}}(I),
\]
the section algebra of the pointwise spatial fibre products, and the map is
the contiguous-block inclusion.  On the endomorphism corner the normalized
trace is
\[
 \tau_{\alpha,q}=P_\alpha(q)\tau_{n,q}.
\]
The germ category at one identifies two morphism sections when they agree on
some smaller interval $[1-\varepsilon,1+\varepsilon]$.  A germ is positive
when it has a positive representative on one such interval.
\end{definition}
\begin{scope}
The ordinary tensor product over $\mathbb C$ is not the section tensor: for
nonconstant $f$, the nonzero element $f\otimes1-1\otimes f$ is killed by
restriction to the diagonal.  Only evaluation at one is canonical on a germ,
and no C*-norm is assigned to germ Hom spaces.
\end{scope}
\provenance{D1141, D1203, D1204}{theory/sidequests/f1-limit/continuous-corners.md}{theory/checks/f1\_operational\_check.py (H3,H13); theory/checks/f1\_limit\_check.py (L4,L14)}{theory/verdicts/f1-limit-adjudication.md}

\begin{theorem}[Continuous corners and their endpoint germ]
{\normalfont\statusProved}\quad
For every compact positive interval $I$, the preceding corners form an
ordered monoidal C*-category.  Their normalized traces are faithful and
multiplicative under ordered tensor; the corner inclusions and their
trace-preserving expectations are continuous and obey the three-block tower
laws.  Fibre evaluation is full.  Germs at one form an ordered monoidal
*-category with a full canonical endpoint evaluation to the $q=1$ corner
category.
\end{theorem}
\begin{scope}
No block exchange is supplied for $q\neq1$.  A nonzero germ may vanish at
the endpoint, which is why endpoint norm is not used as a C*-norm.
\end{scope}
\provenance{F1-LIM-CORNER, F1-LIM-GERM}{theory/sidequests/f1-limit/continuous-corners.md}{H3,H13; L4,L14 on finite fibres and representatives}{theory/verdicts/f1-limit-adjudication.md}
\begin{proof}
Writing $x_\alpha=\sum_{w\in W_\alpha}T_w$, the Hecke multiplication rule
gives $T_ix_\alpha=qx_\alpha$ for every simple reflection in the parabolic.
Thus $x_\alpha^2=P_\alpha x_\alpha$ and $e_\alpha$ is a self-adjoint
projection.  Block length additivity sends
$B_u\otimes B_v$ to the distinct section $B_{u\times v}$, proving
injectivity of the balanced tensor map.  Moreover
$P_{\alpha\mathbin\Vert\beta}=P_\alpha P_\beta$ and
$\iota(e_\alpha\otimes e_\beta)=e_{\alpha\mathbin\Vert\beta}$.
The ambient coefficient expectation is bimodular, so it restricts to each
corner, fixes the included product algebra and preserves the normalized
trace.  All operations commute with restriction and evaluation.  To lift a
fibre corner morphism, lift it with constant normalized-basis coefficients
and compress by the source and target projection sections.  This proves
fullness, both on an interval and after passing to endpoint germs.
\end{proof}

\begin{definition}[Arithmetic and endpoint evaluations]
An arithmetic fibre evaluates the continuous construction at the positive
real number $q=Q=p^r$ and may then apply the complete- or partial-flag
realization over $\mathbb F_Q$.  Endpoint specialization evaluates at
$q=1$.  Coordinate-apartment compression at a fixed arithmetic fibre is a
separate UCP map and is not an evaluation map.
\end{definition}
\begin{scope}
No sequence of prime powers tends to one.  The construction retains the
chosen flag context and does not specialize the full Weyl system or all
polarizations.
\end{scope}
\provenance{D1216}{theory/sidequests/f1-limit/continuous-corners.md}{theory/checks/f1\_operational\_check.py (H9,H13)}{theory/verdicts/f1-limit-adjudication.md}

\begin{proposition}[Arithmetic partial-flag realization]
{\normalfont\statusProved}\quad
For every prime power $Q$ and every composition $\alpha$ of $n$, the fibre
corner $e_\alpha H_n(Q)e_\alpha$ is the full $GL_n(\mathbb F_Q)$-commutant
on partial flags of type $\alpha$, and $P_\alpha(Q)\tau_{n,Q}$ is normalized
physical operator trace.  These identifications preserve corner morphisms,
composition and dagger.
\end{proposition}
\begin{scope}
This is an evaluation followed by a finite-field realization, rather than a
limit over finite fields.  Apartment compression remains nonmultiplicative
when $Q>1$.
\end{scope}
\provenance{F1-LIM-ARITH}{theory/sidequests/f1-limit/continuous-corners.md}{theory/checks/f1\_operational\_check.py (H9,H13)}{theory/verdicts/f1-limit-adjudication.md}
\begin{proof}
A partial flag of type $\alpha$ has $P_\alpha(Q)$ complete refinements.  The
normalized fibre-sum isometry $J_\alpha$ has range projection $e_\alpha$.
Consequently $F\mapsto J_\beta FJ_\alpha^*$ identifies every partial-flag
intertwiner with $e_\beta H_n(Q)e_\alpha$, with inverse
$x\mapsto J_\beta^*xJ_\alpha$.  Ordinary trace on the range of $J_\alpha$
and the flag-count ratio give the asserted normalized trace.
\end{proof}

\subsection{All finite self-adjoint summands}

\begin{definition}[Finite graded self-adjoint completion]
For a nonempty compact positive interval $I$, put
$A_n(I)=H_n^{\mathrm{cts}}(I)$.  An object $X$ is a finite-support family
$(r_n,p_{X,n})_{n\geq0}$, where $r_n\geq0$ and
$p_{X,n}\in M_{r_n}(A_n(I))$ is a self-adjoint projection.  For
$Y=(s_n,p_{Y,n})$, define
\[
 \operatorname{Hom}(X,Y)=
 \bigoplus_n p_{Y,n}M_{s_n,r_n}(A_n(I))p_{X,n}.
\]
Composition and dagger are degreewise matrix multiplication and adjoint,
the identity is $(p_{X,n})_n$, and the norm is the maximum of the finitely
many supremum C*-norms.  Direct sum uses block-diagonal projections within
each degree, so its endomorphism algebra retains the off-diagonal rectangular
corners between equal-degree summands.  Different degrees have zero Homs.
The fibre version at $q$ is denoted $\mathcal U_q^\dagger$.
\end{definition}
\begin{scope}
This is the finite additive dagger Karoubi completion.  It does not assume
duals, rigidity, a spherical trace or finitely many simple objects.
\end{scope}
\provenance{D1261}{theory/sidequests/f1-limit/karoubi-trace.md}{theory/checks/f1\_limit\_karoubi\_check.py (K1)}{theory/verdicts/f1-limit-adjudication.md}

\begin{remark}
An additive object is a coherent multiplicity resource, not a classical
alternative.  For example, the direct sum of two tensor units has
endomorphism algebra $M_2(\mathbb C)$, including off-diagonal coherences.  A
two-outcome classical wire instead has the commutative algebra
$\mathbb C^2$.  Adding finite sums therefore does not change the original
generator tower $\operatorname{End}(x^{\otimes n})=H_n(q)$, and it does not
give the bare one-constituent object qubit observables.
\end{remark}

\begin{definition}[Balanced ordered tensor and finite trace]
For objects $X,Y$, the degree-$k$ multiplicity of $X\otimes Y$ is
\[
 t_k=\sum_{m+n=k}r_ms_n.
\]
Its projection is block diagonal, indexed by $(m,n)$ and lexicographic
matrix indices, with block
\[
 \iota_{m,n}^{(r_m,s_n)}(p_{X,m}\otimes p_{Y,n}).
\]
Morphism tensor is the analogous block-diagonal amplification.  The full
endomorphism corner in $M_{t_k}(A_k(I))$ includes off-diagonal blocks between
different pairs with the same sum.  The unit is the degree-zero object
$(1,1)$, and associators are the permutation matrices identifying the two
triple indexings.  Section tensor is balanced over $C(I)$.

For $a\in\operatorname{End}(X)$ put
\[
 \theta_X(a)=\sum_n(\operatorname{Tr}_{r_n}\otimes\tau_{n})(a_n),
 \qquad w_X=\theta_X(1_X),\qquad \tau_X=\theta_X/w_X.
\]
Here matrix trace is unnormalized.  Operational systems are nonzero objects;
a density satisfies $\rho\geq0$ and $\tau_X(\rho)=1$.
\end{definition}
\begin{scope}
Associators reorder bookkeeping indices but never exchange ordered Hecke
blocks.  The trace is a finite graded matrix/coefficient trace and is not
claimed to arise from categorical duality.
\end{scope}
\provenance{D1262, D1263}{theory/sidequests/f1-limit/karoubi-trace.md}{theory/checks/f1\_limit\_karoubi\_check.py (K1,K2)}{theory/verdicts/f1-limit-adjudication.md}

\begin{theorem}[Full continuous completion and its trace]
{\normalfont\statusProved}\quad
For every nonempty compact positive interval, the finite projection objects
above form an additive ordered monoidal C*-category with every self-adjoint
idempotent split.  The additive self-adjoint completion of the parabolic
corner category is unitarily strong-monoidally equivalent to it.  The trace
$\theta$ is continuous, positive, faithful and cyclic across rectangular
Homs; on a connected interval every nonzero object has $w_X(q)>0$, and
\[
 w_{X\otimes Y}=w_Xw_Y,\qquad
 \tau_{X\otimes Y}\bigl(j(a\otimes b)\bigr)=\tau_X(a)\tau_Y(b).
\]
At a fibre this is the self-adjoint model of the full finite Hecke composition
category.
\end{theorem}
\begin{scope}
The dagger statement concerns self-adjoint projection presentations.  An
algebraic idempotent is merely isomorphic, not canonically unitarily equal,
to its support projection.
\end{scope}
\provenance{F1-LIM-KAR, F1-LIM-KTRACE}{theory/sidequests/f1-limit/karoubi-trace.md}{theory/checks/f1\_limit\_karoubi\_check.py (K1,K2); uniform assertions by proof}{theory/verdicts/f1-limit-adjudication.md}
\begin{proof}
Each Hom is a closed rectangular corner of a finite direct sum of matrix
C*-algebras.  Block diagonal sums give biproducts, and a self-adjoint
idempotent is itself a new projection object.  The normalized block basis
proves injectivity of the balanced tensor, while permutation matrices on
triple indices give the associator and pentagon.  The complete parabolic
composition has projection $1$ in every degree, so its sums and retracts
generate all objects.

For a rectangular $a$, direct expansion gives
\[
 (\operatorname{Tr}\otimes\tau_n)(a^*a)
   =\sum_{i,j}\tau_n(a_{ij}^*a_{ij}),
\]
which proves faithfulness.  Cyclicity follows entrywise.  A continuous
projection has constant ordinary rank on a connected interval, so a nonzero
object cannot disappear at one fibre and $w_X$ stays positive.  Finally the
matrix trace and coefficient trace both multiply on each $(m,n)$ block;
summing over all pairs proves both displayed tensor identities.
\end{proof}

\begin{definition}[Traced assembly of completed objects]
For nonzero $X,Y$, let $B_X=\operatorname{End}(X)$ and let
\[
 j_{X,Y}:B_X\otimes B_Y\longrightarrow B_{X\otimes Y}
\]
be morphism tensor; on intervals its domain is $C(I)$-balanced.  Define
$E_{X,Y}:B_{X\otimes Y}\to B_X\otimes B_Y$ by
\[
 (\tau_X\otimes\tau_Y)(b^*E_{X,Y}(a))
   =\tau_{X\otimes Y}(j_{X,Y}(b)^*a)
\]
for all $b$.  Assembly is the process from the separated word $[X][Y]$ to
the collective atom $[X\otimes Y]$ represented in Heisenberg orientation by
$E_{X,Y}$; split has the reverse type and is represented by $j_{X,Y}$.
\end{definition}
\begin{scope}
The separated algebra and the collective endomorphism algebra are different.
The trace is the finite graded trace just defined, not a pivotal trace.
\end{scope}
\provenance{D1264}{theory/sidequests/f1-limit/karoubi-trace.md}{theory/checks/f1\_limit\_karoubi\_check.py (K1,K5)}{theory/verdicts/f1-limit-adjudication.md}

\begin{theorem}[Continuous traced expectations]
{\normalfont\statusProved}\quad
For nonzero completed objects on a connected compact positive interval,
$j_{X,Y}$ is an injective unital trace-preserving *-map and $E_{X,Y}$ is the
unique normalized-trace-preserving UCP retraction.  These expectations act
continuously on sections, satisfy $Ej=\mathrm{id}$ and the ordered
three-object tower laws, while $jE$ is generally a proper conditional
expectation.  On parabolic objects they recover the earlier corner maps.
\end{theorem}
\begin{scope}
No rigidity or fusion hypothesis is used.  Properness means that collective
same-degree off-diagonal observables can be erased by $jE$.
\end{scope}
\provenance{F1-LIM-KEXPECT}{theory/sidequests/f1-limit/karoubi-trace.md}{theory/checks/f1\_limit\_karoubi\_check.py (K1,K5); complete positivity and continuity by proof}{theory/verdicts/f1-limit-adjudication.md}
\begin{proof}
In one fibre, orthogonal projection onto $j(B_X\otimes B_Y)$ for the faithful
trace inner product fixes the subalgebra and is bimodular.  If the image of a
positive element had a negative spectral projection in the subalgebra, the
trace-pairing identity would be simultaneously negative and nonnegative.
Applying the same argument to matrix amplifications proves complete
positivity.  A local continuous basis turns the trace equations into
$G(q)c(q)=t(q)$ with $G(q)>0$; hence $c(q)=G(q)^{-1}t(q)$ is continuous.
Uniqueness proves agreement on overlaps and the three-object tower law.
\end{proof}
